\documentclass[a4paper,11pt]{article}
\usepackage[scale={0.8,0.9},centering,includeheadfoot]{geometry}
\usepackage{amstext}
\usepackage{amsmath}
\usepackage{verbatim}

\begin{document}
\section*{The ordered Beta-distribution}

\subsection*{Parametrisation}

Let us first summarize the Beta-distribution. It has the following density
\begin{displaymath}
    p(y) = \frac{1}{B(a, b)} y^{a-1}(1-y)^{b-1}, \qquad 0 < y < 1,
    \quad a>0, \quad b > 0
\end{displaymath}
where $B(a, b)$
is the Beta-function 
\begin{displaymath}
    B(a, b) = \frac{\Gamma(a)\Gamma(b)}{\Gamma(a+b)}
\end{displaymath}
and $\Gamma(x)$ is the Gamma-function.
The (re-)parameterisation used is
\begin{displaymath}
    \mu = \frac{a}{a+b}, \qquad 0 < \mu < 1
\end{displaymath}
and
\begin{displaymath}
    \phi = a + b, \qquad \phi > 0,
\end{displaymath}
as it makes
\begin{displaymath}
    \text{E}(y) = \mu \quad\text{and}\quad \text{Var}(y) = \frac{\mu(1-\mu)}{1+\phi}.
\end{displaymath}
The parameter $\phi$ is know as the \emph{precision parameter}, since
for fixed $\mu$, the larger $\phi$ the smaller the variance of $y$.
The parameters $\{a, b\}$ are given as $\{\mu, \phi\}$ as follows,
\begin{displaymath}
    a = \mu \phi \quad\text{and}\quad  b = -\mu \phi + \phi.
\end{displaymath}
The \textbf{ordered Beta}-likelihood (reference
here\footnote{\texttt{https://doi.org/10.1017/pan.2022.20}}) is a
modification which include a derived model for $y=0$ and $y=1$. Let
$\mu = g(\eta)$, where $\eta$ is the linear predictor (see below) and
$g()$ is the inverse link-function, then
\begin{displaymath}
    p_o(y) = \left\{\begin{array}{ll}
      1-g(\eta - k_1) & \text{for } y=0\\
      (g(\eta-k_1) - g(\eta-k_2)) p(y) & \text{for } 0 < y < 1\\
      g(\eta - k_2) & \text{for } y=1
        \end{array}\right.
\end{displaymath}
with two offset-parameters $(k_1, k_2)$ where $k_1 < k_2$.

\subsection*{Link-function}

The linear predictor $\eta$ is linked to the mean $\mu$ using a
default logit-link
\begin{displaymath}
    \mu = \frac{1}{1+\exp(-\eta)}.
\end{displaymath}

\subsection*{Hyperparameter}

The hyperparameter are the precision parameter $\phi$, which is
represented as
\begin{displaymath}
    \phi = s_i \exp(\theta_1)
\end{displaymath}
where $s = (s_i) > 0$ is a fixed scaling. The offset parameters are expressed as
\begin{displaymath}
    k_1 = \theta_1 - \exp(\theta_2), \qquad
    k_2 = \theta_1 + \exp(\theta_2)
\end{displaymath}
The prior is defined on $(\theta_1, \theta_2, \theta_3)$.


\subsection*{Specification}

\begin{itemize}
\item \texttt{family="obeta"}
\item Required argument: $y$
\item Optional argument: $s$ (argument \texttt{scale}, default all 1, $s>0$)
\end{itemize}

\subsubsection*{Hyperparameter spesification and default values}
%% DO NOT EDIT!
%% This file is generated automatically from models.R
\begin{description}
	\item[doc] \verb!The ordered Beta likelihood!
	\item[hyper]\ 
	 \begin{description}
	 	\item[theta1]\ 
	 	 \begin{description}
	 	 	\item[hyperid] \verb!61101!
	 	 	\item[name] \verb!precision parameter!
	 	 	\item[short.name] \verb!phi!
	 	 	\item[output.name] \verb!precision-parameter for the obeta observations!
	 	 	\item[output.name.intern] \verb!intern precision-parameter for the obeta observations!
	 	 	\item[initial] \verb!2.30258509299405!
	 	 	\item[fixed] \verb!FALSE!
	 	 	\item[prior] \verb!loggamma!
	 	 	\item[param] \verb!1 0.1!
	 	 	\item[to.theta] \verb!function(x) log(x)!
	 	 	\item[from.theta] \verb!function(x) exp(x)!
	 	 \end{description}
	 	\item[theta2]\ 
	 	 \begin{description}
	 	 	\item[hyperid] \verb!61102!
	 	 	\item[name] \verb!offset location!
	 	 	\item[short.name] \verb!loc!
	 	 	\item[output.name] \verb!offset location-parameter for the obeta observations!
	 	 	\item[output.name] \verb!offset location-parameter for the obeta observations!
	 	 	\item[initial] \verb!0!
	 	 	\item[fixed] \verb!FALSE!
	 	 	\item[prior] \verb!normal!
	 	 	\item[param] \verb!0 10!
	 	 	\item[to.theta] \verb!function(x) x!
	 	 	\item[from.theta] \verb!function(x) x!
	 	 \end{description}
	 	\item[theta3]\ 
	 	 \begin{description}
	 	 	\item[hyperid] \verb!61103!
	 	 	\item[name] \verb!offset width!
	 	 	\item[short.name] \verb!width!
	 	 	\item[output.name] \verb!offset width-parameter for the obeta observations!
	 	 	\item[output.name] \verb!offset width-parameter for the obeta observations!
	 	 	\item[initial] \verb!0!
	 	 	\item[fixed] \verb!FALSE!
	 	 	\item[prior] \verb!loggamma!
	 	 	\item[param] \verb!1 1!
	 	 	\item[to.theta] \verb!function(x) log(x)!
	 	 	\item[from.theta] \verb!function(x) exp(x)!
	 	 \end{description}
	 \end{description}
	\item[status] \verb!experimental!
	\item[survival] \verb!FALSE!
	\item[discrete] \verb!FALSE!
	\item[link] \verb!default logit loga cauchit probit cloglog ccloglog loglog!
	\item[pdf] \verb!obeta!
\end{description}


\subsection*{Example}

In the following example we estimate the parameters in a simulated
example.
\verbatiminput{example-obeta.R}


\subsection*{Notes}

None.

\end{document}


% LocalWords:  hyperparameter overdispersion Hyperparameters nbinomial

%%% Local Variables: 
%%% TeX-master: t
%%% End: 
